% !TEX root = ../main.tex

\begin{acknowledgements}
  朝来庭下，光阴似箭。于上海交通大学电子系求学八载，始于少年意气的憧憬，终于此际而立之年落笔时的沉静。炽热的青春与深井般的科研困顿，赋予论文以独特的厚度。信手翻去，某页的缘起，或许是19年秋日在402实验室的灵光乍现，而下一章的框架，可能来源于23年在3507实验室的一次组会讨论。然回首八载岁月，科研之径从来荆棘多于坦途。其间所历，是面对文献浩海不得其门时的长久彷徨，是一次次自以为是的构思被证伪时的沮丧，是目睹朋友在各个赛道高歌猛进时自身课题停滞不前的自我否定，是对家人报喜不报忧的小心谨慎。我的记忆也跟着明灭起伏不定起来。对暗处的记忆是模糊的，彼时是对科研纯粹性崩塌后的虚无，是科研停止不前对自我能力的解构。然而在这片晦暗的记忆旷野中，总有一些明媚的瞬间，是老师，是家人，是同门，是朋友。是他们，在我闭门造车时，伸手将我拉出了那间停滞不前的尺椽之室。在明灭盘桓之间，我开始模糊地触摸到科研中那组“绝对”与“相对”的张力：所谓绝对，是科研目标与研究空白领域如重力般客观存在；所谓相对，是科研目标的达成从来不是一蹴而就的，需于约束中寻优，在迭代中趋近从而寻找破局之点。倘若将科研的信念、自身的努力、长年的专注、与师友亲人的托举视作一张多维的支撑之网，那么其中任何一维的松脱都足以令我滑落。我深知，就是因为这张大网最终让我得以带着几分踉跄，在这博士学位的隘口站稳并得以入门学术殿堂。若要我为自己这八年作一概括，我想起的便是那句：“我用尽了全力，过着平凡的一生。” 我庆幸这全力，为我赢得一份继续进行学术探索的资格。就是这张由无数善意、帮助与自身坚持编织的网，让我得以体面地踏入学术之门。此刻，我郑重地表达对所有帮助我、宽慰我、支持我的人的感谢：

  感谢我的导师归琳教授。她是我科研道路上的引路人，以独特的视野为我打开了工业界的大门，使我得以洞悉真实的行业痛点与产业链逻辑。归老师常言：“要清楚自己的研究处在什么位置。”这句话如座右铭般刻在我心中，使我始终自觉追问：一项工作的源头何在？它究竟回应了哪个层面的问题？又将如何作用于真实的物理世界？其次，归老师赋予我的科研以极大的自由。她鼓励我跟随兴趣自主探索，这份信任极大地点燃了我的内在动力；而每当我思路行将偏斜，她又总能于关键处给予精准的提点，让我不至迷失。我也深深感谢归老师对我参与国际交流的坚定支持。那些与不同文化背景的学者碰撞思想的经历，极大地拓宽了我的学术视野。最后，于我而言，归老师更是一位人生导师。她以身作则，教我如何以严谨务实的态度工作，以开放有效的方式沟通；她身上那种锐意进取、敢于冒险、勇于挑战的精神气质，已悄然内化为我前行的一部分。

  感谢陈昊鹏教授。陈老师是我学术研究的启蒙者。他总能精准地捕捉到我脑海中那些粗浅的想法，并以其清晰的逻辑与严谨的框架，帮助我将它们转化为可落地、可验证的技术方案，最终凝结为具学术创新性的论文。每当我写作陷入停滞，陈老师总能以其细致的阶段性规划与精辟的针对性提点，为我拨开迷雾并推进论文进展。此外，陈老师让我深切体认到团队协作的力量。他曾主动将研究方向相近的同学们组织起来形成研讨小组。在那种开放的协作、高频的讨论与技术的取长补短中，我不仅看到了技术如何在他人的激发下迅速迭代，更感受到了集体智慧如何把每个人的论文与观点打磨得更具创新性。

  感谢项目组的俞晖老师、任锐老师与应凯老师。首先，由衷感谢他们的信任，让我有幸参与到重点、重大工程项目之中，学习如何将所学知识应用在工程实际，真正将“论文写在祖国大地”。其次，我要深深感谢三位老师对我工程思维的培养。工程思维迥异于纯粹的学术思维，它要求算法在复杂的系统约束与多维的耦合关系中寻求可行解。正是老师们一次次耐心的指正与谆谆教导，让我逐步养成了从系统全局出发、在约束中寻找权衡的思维习惯。最后，特别感谢老师们对我实践能力的悉心培养。他们让我深刻认识到，若只停留在提出算法与理论优化，无异于纸上谈兵。在真实的项目历练中，我不仅锻炼了用代码严谨验证想法的能力，更掌握了在复杂现实系统中将理论落地的关键技能。这段经历，是我从学术构想走向工程实现的关键桥梁。
  
  感谢在我实验系统搭建过程中倾力相助的师弟师妹们：孙琳、黄沛琪、高肖健、诸柏铭、刘津铭。你们的支持不仅为我节省了宝贵的时间，更让我得以在严谨、高效的验证中，不断优化与迭代方案。

  感谢朝夕相处、并肩奋斗的3507实验室全体同门：孙飞、张琦、秦启波、朱世超、张凌、刘畅、莫潇豪、桑锡超、邓慧敏、赵构恒、马蕾、张勇成、徐培凯、肖迪笙、吴泽、李昊、魏麟懿、程妍璇、胡嘉弘。那些在组会中与你们进行的深度讨论与思维碰撞打磨了我的论文。能与这样一群优秀而温暖的同行者共度这段时光，是我莫大的幸运。

  最后，由衷地感谢我的父母。你们不仅给予我生命，让我得以历经人生的山峰与河流，更为我备好行囊、筑起港湾。一路走来，无条件的支持、温暖的陪伴与无尽包容的精神鼓励，是你们最深沉的爱。你们的付出与无私，让我得以心无旁骛地行走至今。特别感谢我的女朋友何文轩女士。你的到来，合上了我人生中至关重要的一块拼图。从此，我的生命拥有了最坚实的地基，静待着我们共同的未来生根发芽。然而这份等待漫长而煎熬，感谢你始终以爱与包容，为我扫除阴霾，注入信心。你总是鼓励我、认可我，正是这份坚定的信任，最终赋予了我百折不挠的勇气与一往无前的认真。

  文末搁笔，心绪翻涌。这八年的读研经历所赋予我的，远不止一纸文凭，更是对我意志品性与社会能力的改造与重塑。我深知，迈出校园仅是人生真正的序章，未来的挑战或许不亚于读博的艰辛。但此刻，我心中却充满前所未有的平静。我忽然懂得，琢之磨之，玉汝于成。这八年的读博经历，是一场漫长而专注的自我琢磨。往后的岁月，亦当如是。以生活为砧，以经历为砺，不避寒暑，不惧切磋，将生活作为一件最郑重的作品，从容不迫地，继续雕琢下去。
  
  言辞有尽，谢意难全。最后向所有关心我的家人、师长和朋友们表示深深的谢意，并向百忙之中参加论文评阅和答辩的诸位老师致以诚挚的谢意！
\end{acknowledgements}
